\uuid{EO3D}
\titre{Équation droite/plan}
\auteur{bodin ; hanani ; mzari}
\organisation{exo7}
\licence{CC BY-NC-SA 4.0}
\datecreate{2020-01-20}
\niveau{L1}
\module{Géométrie}
\chapitre{Géométrie du plan}
\sousChapitre{Distance, norme, produit scalaire}
\difficulte{3}
\theme{équation droite, plan}
\exo7qcmid{lille-1.geometrie-plane.024}

\contenu{
\begin{question}
\qclassification{140.01,140.02}
\qtags{motcle=équation droite/plan}
On considère trois points $A$, $B$ et $C$ du plan tels que
$$(AB)\; : \; x-2y+3=0\quad \mbox{ et }\quad (AC)\; : \; 2x-y-3=0.$$
Quelles sont les bonnes réponses ?
\begin{answers}
    \bad{Les points $A$, $B$ et $C$ sont alignés.}
    \bad{Le point $B$ appartient à $(AC)$.}
    \bad{Le point $C$ appartient à $(AB)$.}
    \good{Les coordonnées de $A$ sont $A(3,3)$.}
\end{answers}
\begin{explanations}
Les points $A$, $B$ et $C$ ne sont pas alignés car sinon les droites $(AB)$ et $(AC)$ seraient confondues. Ces droites se coupent en $A$ et les coordonnées de ce point d'intersection vérifient le système
$$\left\{\begin{array}{l}x-2y+3=0\\ 2x-y-3=0
\end{array}\right. \Leftrightarrow \left\{\begin{array}{l}x=3\\ y=3.\end{array}\right.$$
\end{explanations}
\end{question}
}
